\section*{PTT finite temporal witness synchronization lemma}

This lemma discharges only the finite synchronization step following the
temporal witness survival theorem interface.

\paragraph{Hypotheses.}
Fix a derivation \(D\), a finite nonempty set \(\mathcal W\) of witnesses,
and a predicate
\[
\mathrm{TemporalConfigurationWitness}(D,t,W).
\]
Assume
\[
\forall W\in\mathcal W,\;
\exists t\in\mathbb N,\;
\mathrm{TemporalConfigurationWitness}(D,t,W).
\]

\paragraph{Lemma.}
There exists \(T\in\mathbb N\) such that
\[
\forall W\in\mathcal W,\;
\exists t\leq T,\;
\mathrm{TemporalConfigurationWitness}(D,t,W).
\]

Moreover, there exist \(W_\ast\in\mathcal W\) and \(t_\ast\leq T\) such that
\[
\mathrm{TemporalConfigurationWitness}(D,t_\ast,W_\ast).
\]

\paragraph{Proof.}
For each \(W\in\mathcal W\), define
\[
\tau(W)
=
\min\left\{
t\in\mathbb N:
\mathrm{TemporalConfigurationWitness}(D,t,W)
\right\}.
\]
The defining set is nonempty by hypothesis, so \(\tau(W)\) exists by the
well-ordering principle for \(\mathbb N\).

Since \(\mathcal W\) is finite and nonempty, the finite nonempty set
\[
\{\tau(W):W\in\mathcal W\}
\]
has a maximum. Define
\[
T=\max\{\tau(W):W\in\mathcal W\}.
\]
Then for every \(W\in\mathcal W\),
\[
\tau(W)\leq T
\]
and
\[
\mathrm{TemporalConfigurationWitness}(D,\tau(W),W).
\]
Taking \(t=\tau(W)\) proves
\[
\forall W\in\mathcal W,\;
\exists t\leq T,\;
\mathrm{TemporalConfigurationWitness}(D,t,W).
\]

Because the maximum is attained, there exists \(W_\ast\in\mathcal W\) with
\[
\tau(W_\ast)=T.
\]
Taking \(t_\ast=T\) gives
\[
\mathrm{TemporalConfigurationWitness}(D,t_\ast,W_\ast).
\]
\(\square\)

\paragraph{Boundary.}
This lemma does not construct any
\(\mathrm{TemporalConfigurationWitness}\).
It requires the individual witnesses as hypotheses.
It does not prove a verified trace extraction theorem.
It does not construct \(\mathrm{TemporalWitnessSource}\).
It does not prove PTT closure.
